% Nouveau EGA V (cote 25): archive cote 25, site page 43 (PDF page 44).
% typescript. Draft compared: none.
% AI-transcribed by Claude Opus 5.5 (claude-opus-5-5, adaptive thinking, effort high), Message Batches API, batch ega25b (msgbatch_01A7n5YuXubUPNK1zsepi8RG), 2026-11-13; not yet checked by eye.
\scanpage{25/43}
% NOTE: The dot below the second Y in(ii) has uncertain function. Another isolated dot lies below the opening parenthesis around the final hypothesis number; no additional character value is assigned.
\mspage{\hand{39}}

L'adhérence de son images dans $\overset{\hand{\vee}}{P}$ , muni de la structure réduite
induite, et considérons le morphisme induit g: $\underline{Y}^{\text{sing}} \hand{\longrightarrow} T$ \var{\xout{de}}{}{typewriter strike-out}{ega25_25_043-v0001} (morphis-%
me dominant de préschémas \var{\xout{irréductibles}}{}{typewriter strike-out}{ega25_25_043-v0002} intègres). Alors les condition
\var{\inserted[typed]{(i) à (ii bis)}}{(i) à (ii bis)}{insertion}{ega25_25_043-v0003} suivantes sont équivalentes:

(i) Le morphisme g est génériquement \var{\deleted{non ramifié}}{}{deletion}{ega25_25_043-v0004} \var{\inserted{étale}}{étale}{insertion}{ega25_25_043-v0005} (i.e. est \var{\deleted{gon
ramifié}}{}{deletion}{ega25_25_043-v0006} \var{\inserted{étale}}{étale}{insertion}{ega25_25_043-v0007} en au moins un point, ou ce qui revient au même, est \var{\inserted{étale =}}{étale =}{insertion}{ega25_25_043-v0008} non ramifié
au point générique de $\underline{Y}^{\text{sing}}$).

(ibis) L'extension de corps L/K définie par g est finie et
séparable.

\var{\margin{NB. L'hyp. f une \uline{imm.} (au lieu de : non ramifié) ne sert que pour l'équivalence avec (i ter). Énoncer 8.8. en conséquence\dots}}{\placedmargin{NB. L'hyp. f une \uline{imm.} (au lieu de : non ramifié) ne sert que pour l'équivalence avec (i ter). Énoncer 8.8. en conséquence\dots}}{marginal note placed}{ega25_25_043-v0009}
(i ter) Le m\unsure{oo|or}phisme g est birationnel, i.e. l'extension L/K est
l'extension triviale.

(ii) $\underline{Y}^{\text{sihg}} \hand{\neq} \underset{\hand{\unsure{\cdot}}}{\underline{Y}}^{\text{sup sing}}$ \var{\inserted{(ensemblistement, disons).}}{(ensemblistement, disons).}{insertion}{ega25_25_043-v0010}

(ii bis) Il existe un $x \hand{\in} X(\bar{k})$ , et un \var{\over{h}{h}}{h}{writing over}{ega25_25_043-v0011}yperplan tangent \var{\inserted[typed]{H}}{H}{insertion}{ega25_25_043-v0012} à $X_{\bar{k}}$
en x qui ne soit pas ``osculateur'' en x (par quoi on entenda précisément
que x n'est pas supersingulier pour la section de $\underline{O}_{X_{\bar{k}}}(1)$ qui définit
H ... ).

Ces conditions impliquent que \var{\inserted{$\underline{Y}^{\text{sup sing}} \neq \emptyset$ ou}}{$\underline{Y}^{\text{sup sing}} \neq \emptyset$ ou}{insertion}{ega25_25_043-v0013} dim $\underline{Y}^{\text{sup sing}}$ \var{\deleted{$\leq$}}{}{deletion}{ega25_25_043-v0014}$= r-2$, donc que
l'image de $\underline{Y}^{\text{sup sing}}$ dans $\overset{\hand{\vee}}{P}$ a une dodimension $\hand{\geqslant} 2$, et \var{\xout{que T
est de dimension r-1, i.e. de codimension 1 dans $\overset{\hand{\vee}}{P}$}}{}{typewriter strike-out}{ega25_25_043-v0015} elles impliquent
également:

(iii) dim T = r-1, i.e. T est de codimension 1 dans $\overset{\hand{\vee}}{P}$ .

Dém. L'équivalence de (i) et (i bis) est triviale , son équiva-%
lence avec (ii) est une conséquence triviale de 8.\var{\over{\unsure{5|6}}{7}}{7}{writing over}{ega25_25_043-v0016}. b), enfin l'équi-%
valence de (ii) et (ii bis) est pratiquement la définition de
$\underline{Y}^{\text{sup sing}}$ . Evidemment (i ter) \hand{$\Longrightarrow$}(i) , il reste à prouver que (i)\hand{$\Longrightarrow$}
(i ter). On peut évidemment supposer k alg clos, et on est ramené à
prouver (compte tenu de l'hypothèse (\unsure{ii|iii}) ) qu\var{\xout{e}}{}{typewriter strike-out}{ega25_25_043-v0017}'il existe un ouvert U\hand{$\neq\emptyset$}

\var{\margin{\hand{43}}}{\placedmargin{\hand{43}}}{marginal note placed}{ega25_25_043-v0018}
